\documentclass[10pt,a4paper]{amsart}
\usepackage[margin=19mm]{geometry}
\usepackage{amsmath,amssymb,mathtools}
\usepackage[scr=rsfs,cal=euler]{mathalfa}
\usepackage{xfrac}
\usepackage[unicode,hidelinks,bookmarks=false]{hyperref}
\newcommand{\ie}{i.e.\ }
\newcommand{\cf}{cf.\ }
\newcommand{\resp}{respectively\ }
\newcommand{\Subsection}{\subsection}
\providecommand{\nobreakdash}{\nobreak\hbox{-}}
\setlength{\emergencystretch}{2em}
\multlinegap0pt
\usepackage{bm}
\usepackage{bbm}
\usepackage{dsfont}
\usepackage{xspace}
\usepackage{cancel}
\usepackage{mathtools}
\usepackage{amsthm}
\makeatletter
\@ifundefined{theorem}{\newtheorem{theorem}{Theorem}}{}
\@ifundefined{proposition}{\newtheorem{proposition}[theorem]{Proposition}}{}
\makeatother
\makeatletter
\expandafter\ifx\csname proof\endcsname\relax
\newenvironment{proof}[1][Proof]{\noindent\textbf{#1.} }{\ \rule{0.5em}{0.5em}}
\fi
\makeatother
\title[A quantitative stability result for regularity of optimal tr]{A quantitative stability result for regularity of optimal transport on compact manifolds}
\author{Micah Warren}
\date{Source: arXiv:2507.09035v1 (CC BY 4.0)}
\begin{document}
\maketitle

\noindent\emph{Source and licence.} A quantitative stability result for regularity of optimal transport on compact manifolds. Version arXiv:2507.09035v1 (CC BY 4.0), distributed under the Creative Commons Attribution 4.0 International licence, as recorded in the arXiv licence metadata for this exact version and shown on the abstract page https://arxiv.org/abs/2507.09035v1. Full rights notice: licenses/arxiv-2507.09035-CC-BY-4.0.txt.\par\smallskip

\noindent\textit{Editorial scope.} Counted unit: the Proposition whose source label is \texttt{bbbb} in the article (source0.tex lines 524--538, source label \texttt{bbbb}), delivered together with the article's own complete original proof of it (source0.tex lines 540--612, one \texttt{proof} environment of 73 lines). No mathematics is written, repaired or re-derived: statement and proof are the article's own text line for line. Required context retained uncounted: meq (lines 264--267). Every definition, display and label of those lines is retained unchanged. Omitted from this package and not redistributed: 1--263, 268--523, 539--539, 613--1065. Nothing inside a retained excerpt is omitted or altered: every retained body is delivered exactly as the article's lines are written, so no omission receipt of any kind accompanies this package. Citations: the retained excerpts carry no citation command at all, so no bibliography entry of the article is transcribed and no reference is redistributed. Reference adaptation: none was needed and none was made; every reference inside the delivered unit resolves inside it, so no unresolved reference or citation is produced. Printed slips kept literal: no printed word, symbol, hypothesis or number of the counted statement or of its proof was altered. Layout and class: the article's own class is \texttt{amsart} with packages including amsmath, amsfonts, amssymb, graphicx; the wrapper is the lane's pinned \texttt{amsart} editorial wrapper at 10pt on A4, which restates the theorem-like environments the retained text needs and adds the article's own macro definitions that the retained text needs (none), with extra wrapper packages (bm, bbm, dsfont, xspace, cancel, mathtools). The delivered numbering is the unit's own. Assets: no retained span contains a figure environment, a graphics-inclusion command, a PDF-page inclusion, a table or a TikZ picture; no image file is copied into this package, the package declares no asset dependency and no image file is redistributed with it. Licence: version arXiv:2507.09035v1 of this article is distributed under the Creative Commons Attribution 4.0 International licence, as recorded in the arXiv licence metadata for this exact version and shown on the abstract page https://arxiv.org/abs/2507.09035v1. A full rights notice is delivered at licenses/arxiv-2507.09035-CC-BY-4.0.txt. Characters: every retained excerpt is pure ASCII, so the delivered engine, XeLaTeX with its default Latin Modern text font, reports no missing character.
\begin{equation}
0.9\left\Vert \xi\right\Vert _{g_{0}}^{2}\leq\sum_{i=1}^{n}\left(  \xi
^{i}\right)  ^{2}\leq1.1\left\Vert \xi\right\Vert _{g_{0}}^{2} \label{meq}.%
\end{equation}

\begin{proposition}
\label{bbbb}Suppose that $u$ is a semi-convex function on $M$. \ Let
\[
b=\max\left\vert \nabla u\right\vert .
\]
Then there is a value $c(A,g,\mu)>0$ and a value $C(A,g,\mu)$ such that if
\[
\int\left\vert u\right\vert ^{2}d\mu\leq c(A,g,\mu)
\]
then
\[
b^{n+4}\leq C(A,g,\mu)\left\Vert u\right\Vert _{L^{2}(d\mu)}^{2}.
\]

\end{proposition}

\begin{proof}
Take normal coordinates at the point where $\left\vert \nabla u\right\vert $
is maximized, such that
\[
\frac{\nabla u}{\left\vert \nabla u\right\vert }=\partial_{1}.%
\]
There is a ball of radius $1$ such that on $B_{1}$ (recalling (\ref{meq})) that
\begin{equation}
\left\vert Du\right\vert \leq1.1b \label{meq2}.%
\end{equation}
Also, using
\[
u_{11}\geq-A
\]
we have
\[
u_{1}(x_{1},0,...,0)\geq b-Ax_{1}\text{ for }x_{1}>0.
\]
Now first we consider the case that
\[
b<A.
\]
Then
\begin{align*}
u(\frac{b}{2A},...,0)-u(0)  &  \geq\int_{0}^{\frac{b}{2A}}\left(
b-Ax_{1}\right)  dx_{1}\\
&  =\frac{3b^{2}}{8A}.
\end{align*}
Now let
\begin{align*}
Q_{1}  &  =B_{\frac{b}{11A}}(0)\\
Q_{2}  &  =B_{\frac{b}{11A}}(\frac{b}{2A},...,0))
\end{align*}
Then by (\ref{meq2})
\begin{align*}
\max_{x\in Q_{1}}\left(  u(x)-u(0)\right)   &  \leq1.1b\frac{b}{11A}\\
\min_{x\in Q_{1}}\left(  u(x)-u(\frac{b}{2A},...,0)\right)   &  \geq
-1.1b\frac{b}{11A}.%
\end{align*}
Thus
\begin{align*}
\max_{x\in Q_{1}}u(x)  &  \leq\frac{b^{2}}{10A}+u(0)\\
\min_{x\in Q_{2}}u(x)  &  \geq-\frac{b^{2}}{10A}+u(\frac{b}{2A},...,0)\\
&  \geq-\frac{b^{2}}{10A}+\frac{3b^{2}}{8A}+u(0)\\
&  \geq\frac{7b^{2}}{40A}+\max_{x\in Q_{1}}u(x)
\end{align*}
which means on at least one of the two balls, we have
\[
\left\vert u\right\vert \geq\frac{7b^{2}}{80A}.
\]
It follows that
\begin{align*}
\left\Vert u\right\Vert _{L^{2}(d\mu)}^{2}  &  \geq\min\mu(x)c_{n}\left(
\frac{b}{11A}\right)  ^{n}\left(  \frac{7b^{2}}{80A}\right)  ^{2}\\
&  \geq c(A,\mu)b^{n+4}.
\end{align*}
Next we claim that $b\geq A$ can be excluded by choosing $c(A,g,\mu)$ small
enough. \ This follows a nearly identical argument, replacing
\begin{align*}
&  u(\frac{b}{2A},...,0)\text{ with }u(\frac{1}{2},...,0)\\
&  B_{\frac{b}{11A}}(0)\text{ with }B_{\frac{1}{11}}(0)
\end{align*}
giving%
\[
u(\frac{1}{2},...,0)-u(0)\geq\frac{3b}{8}%
\]
etc, and showing that
\[
\left\Vert u\right\Vert _{L^{2}(d\mu)}^{2}\geq\min\mu(x)c_{n}\left(  \frac
{1}{11}\right)  ^{n}\left(  \frac{7}{80}b\right)  ^{2}.
\]

\end{proof}

\end{document}
